%!TEX root = ../thesis.tex
\newpage
\thispagestyle{empty}
% If you do want an image in the colophon:
\begin{figure}
  \vspace{50pt}
  \centering
    \includegraphics[width=0.50\textwidth]{endmatter/colophon1.png}
\end{figure}

% If you don't want an image in the colophon:
% \vspace*{200pt}

\begin{center}
\parbox{200pt}{\lettrine[lines=3,slope=-2pt,nindent=3pt]{\textcolor{SchoolColor}{Γ}}ια την ανάπτυξη του παρόντος εγγράφου χρησιμοποιήθηκε \LaTeX, μια γλώσσα ανοιχτού τυπογραφικού προγραμματισμού βασισμένη στην \TeX \, του Donald Knuth, και με δημιουργό τον Leslie Lamport. Το πρότυπο που χρησιμοποιήθηκε, και φέρει παρόμοια εμφάνιση και αισθητική, δημιουργήθηκε από τον Krishna Kumar:\href{mailto:kks32@cam.ac.uk}{kks32@cam.ac.uk}, κυκλοφορεί υπό την άδεια \textsc{MIT}(\textsc{x}11), και είναι διαθέσιμο στο σύνδεσμο \url{https://github.com/kks32/phd-thesis-template/}. Το κείμενο συντάσσεται σε γραμματοσειρά μεγέθους 12 μονάδων και φόντου XITS, ένα ελεύθερο φόντο που χρησιμοποιείται σε επιστημονικά έγγραφα. Η παραπάνω εικόνα, \textit{Book Vignette Cityscape}, δημιουργήθηκε από το εργαστήρι \href{http://buysellgraphic.com/vector-graphic/download/background_vector_illustration_with_book_and_vignette_cityscape_26398.html}{BSGStudio} και δεσμεύεται από τους όρους \href{https://creativecommons.org/licenses/by/4.0/deed.en}{\textsc{CC BY-NC-ND 4.0}}.}
\end{center}